\chapter{Template Metaprogramming}
\label{ch:template-metaprogramming}

\textit{Templates are a Turing-complete functional language that runs inside the compiler. C++11 gave it variadics, type traits, and \texttt{constexpr} --- turning a clever trick into an engineering discipline.}

Template metaprogramming (TMP) generates and selects code at \textbf{compile time}. C++11 transformed TMP from arcane to practical with \textbf{variadic templates}, a real \textbf{\texttt{<type\_traits>}} library, \textbf{alias templates}, \textbf{\texttt{constexpr}}, and \textbf{\texttt{std::tuple}}. This chapter develops the functional model of templates, then the toolkit: parameter packs, metafunctions, type traits, SFINAE, the detection idiom, tag dispatch, and tuples.

%% ============================================================
\section{Templates as a Compile-Time Functional Language}

The C++ template system is, by itself, a \textbf{pure functional language evaluated by the compiler}. It has no mutation and no loops; ``variables'' are types and compile-time constants, and the only ``looping'' construct is \textbf{recursion}, terminated by a \textbf{base case} (a template specialization). By convention a \emph{template metafunction} exposes its result through a member named \texttt{::type} (type results) or \texttt{::value} (value results).

\begin{lstlisting}[language=C++, caption={a metafunction computing a value at compile time}]
template<int N>
struct Factorial {
    static const int value = N * Factorial<N - 1>::value;
};
template<>
struct Factorial<0> {            // base case (recursion terminator)
    static const int value = 1;
};
// Factorial<5>::value == 120, computed entirely by the compiler
int buffer[Factorial<5>::value]; // usable as a compile-time constant
\end{lstlisting}

%% ============================================================
\section{Variadic Templates and Parameter Packs}

A \textbf{variadic template} accepts an arbitrary number of template arguments via a \textbf{parameter pack}. Three pieces of syntax govern packs: \texttt{typename... Args} (template parameter pack), \texttt{Args... args} (function parameter pack), and \texttt{sizeof...(Args)} (the number of elements). A pack is consumed by \textbf{expansion} (\texttt{pattern...}), typically via recursion: peel the first argument, recurse on the rest, stop at a non-variadic base case.

\begin{lstlisting}[language=C++, caption={classic recursive variadic print}]
#include <iostream>
void print() {}                                   // base case

template<typename T, typename... Args>
void print(T first, Args... args) {
    std::cout << first << " ";
    print(args...);                               // expand the rest
}
int main() { print(1, "hello", 3.14); }           // 1 hello 3.14
\end{lstlisting}

Packs combine with perfect forwarding (Chapter~\ref{ch:move-semantics-and-smart-pointers}): \texttt{std::forward<Args>(args)...} expands each argument while preserving its value category --- the mechanism behind \texttt{make\_unique}, \texttt{make\_shared}, and \texttt{emplace\_back}.

%% ============================================================
\section{Iterating Over a Parameter Pack}

Recursion is the portable C++11 way to process every element. A second C++11 technique is the \textbf{expander trick}, which performs the whole operation in one function using a brace-initialized dummy array to force left-to-right evaluation:

\begin{lstlisting}[language=C++, caption={the C++11 "expander" trick (no second overload)}]
template <class... Ts>
void print_all(std::ostream& os, Ts const&... args) {
    using expander = int[];
    (void)expander{ 0, (void(os << args), 0)... }; // os << arg for each
}
\end{lstlisting}

\begin{quote}
\textbf{C++14/17 forward references:} C++17 \textbf{fold expressions} reduce this to \texttt{((os << args), ...);}, and \textbf{\texttt{if constexpr}} lets the recursive form live in one function by compiling the recursive call only when \texttt{sizeof...(rest) > 0}. Neither exists in C++11.
\end{quote}

%% ============================================================
\section{Template Metafunctions}

C++11 offers cleaner spellings of the factorial idea. Inheriting from \texttt{std::integral\_constant} gives \texttt{::value}, a \texttt{value\_type}, and a \texttt{constexpr} conversion for free:

\begin{lstlisting}[language=C++, caption={a metafunction via std::integral\_constant}]
#include <type_traits>
template<long long n>
struct factorial : std::integral_constant<long long, n * factorial<n-1>::value> {};
template<>
struct factorial<0> : std::integral_constant<long long, 1> {};
// factorial<7>::value == 5040
\end{lstlisting}

A \texttt{constexpr} function (Chapter~\ref{ch:the-modern-c11-core}) is usually cleanest, and --- unlike a metafunction --- also works at runtime when its arguments are not compile-time constants:

\begin{lstlisting}[language=C++, caption={constexpr serves both compile time and runtime}]
constexpr long long factorial(long long n) {
    return n == 0 ? 1 : n * factorial(n - 1); // C++11: single return
}
template <typename T>
constexpr T power(T value, unsigned exp) {
    return exp == 0 ? 1 : value * power(value, exp - 1);
}
void use() {
    constexpr int compile_time = power(3, 3); // computed by the compiler
    int x; std::cin >> x;
    int run_time = power(x, 3);               // same code, evaluated at runtime
}
\end{lstlisting}

\begin{quote}
\textbf{C++14 forward reference:} C++14 lifts the single-return restriction, so \texttt{constexpr} functions may use \texttt{if}/\texttt{for} directly.
\end{quote}

%% ============================================================
\section{Type Traits}

\texttt{<type\_traits>} provides compile-time \textbf{inspection} (\texttt{::value}) and \textbf{transformation} (\texttt{::type}) of types.

\begin{lstlisting}[language=C++, caption={inspecting types}]
#include <type_traits>
static_assert(std::is_integral<int>::value, "int is integral");
static_assert(std::is_pointer<int*>::value, "int* is a pointer");
static_assert(std::is_same<int, signed int>::value, "same type");
\end{lstlisting}

Writing your own trait follows a three-step pattern: a primary template inheriting \texttt{std::false\_type}, a specialization inheriting \texttt{std::true\_type}, and a front-end that strips qualifiers:

\begin{lstlisting}[language=C++, caption={implementing is\_pointer from scratch}]
template <typename T> struct is_pointer_      : std::false_type {};
template <typename T> struct is_pointer_<T*>  : std::true_type  {};
template <typename T> struct is_pointer
    : is_pointer_<typename std::remove_cv<T>::type> {};   // strip const/volatile
static_assert(is_pointer<int* const>::value, "yes");
\end{lstlisting}

Compile-time \texttt{if}/\texttt{else} on types is \texttt{std::conditional}:

\begin{lstlisting}[language=C++, caption={select a type at compile time}]
template<typename T>
struct ValueOrPointer {
    typename std::conditional<(sizeof(T) > sizeof(void*)), T*, T>::type vop;
};
\end{lstlisting}

%% ============================================================
\section{SFINAE}

\textbf{SFINAE} --- \emph{Substitution Failure Is Not An Error} --- is the rule that makes constrained overloading work: if substituting template arguments produces an ill-formed type or expression \textbf{in the immediate context}, that candidate is silently \emph{removed} from overload resolution rather than causing a hard error.

\begin{lstlisting}[language=C++, caption={SFINAE removes a non-viable candidate}]
template <class T>
auto begin(T& c) -> decltype(c.begin()) { return c.begin(); } // #1
template <class T, std::size_t N>
T* begin(T (&arr)[N]) { return arr; }                         // #2
int vals[10];
begin(vals); // #1 fails (arrays have no .begin()); discarded, #2 selected.
\end{lstlisting}

Only failures \textbf{in the immediate context} are soft; an error deep inside an instantiated body is hard. \texttt{std::enable\_if} is the standard switch: \texttt{enable\_if<true, R>::type} is \texttt{R}; \texttt{enable\_if<false, R>::type} does not exist, so any overload depending on it vanishes.

\begin{lstlisting}[language=C++, caption={constraining an overload with enable\_if}]
template<typename T>
typename std::enable_if<std::is_integral<T>::value, T>::type
gcd(T a, T b) { return b == 0 ? a : gcd(b, a % b); } // integral T only
\end{lstlisting}

%% ============================================================
\section{The Detection Idiom (\texttt{void\_t})}

To ask ``is this expression valid for type \texttt{T}?'' use the \textbf{detection idiom}, built on \textbf{\texttt{void\_t}} --- a metafunction mapping any well-formed type list to \texttt{void}.

\begin{lstlisting}[language=C++, caption={void\_t (a C++11 one-liner; std::void\_t in C++17)}]
template <class...> using void_t = void;
\end{lstlisting}

\begin{lstlisting}[language=C++, caption={detect a member function foo()}]
#include <type_traits>
#include <utility>   // std::declval
template <class T, class = void>
struct has_foo : std::false_type {};
template <class T>
struct has_foo<T, void_t<decltype(std::declval<T&>().foo())>> : std::true_type {};

struct A {};
struct B { void foo(); };
static_assert(!has_foo<A>::value, "A has no foo");
static_assert( has_foo<B>::value, "B has foo");
\end{lstlisting}

When \texttt{T.foo()} is well-formed, the specialization's second argument resolves to \texttt{void} and is preferred by partial ordering; otherwise substitution fails softly and the \texttt{false\_type} primary remains. The same pattern detects operators, nested types, and \texttt{std::hash} specializations.

%% ============================================================
\section{Tag Dispatching}

\textbf{Tag dispatching} selects an overload at compile time by passing an empty \emph{tag} type (often a trait result) as an extra argument --- frequently more readable than nested \texttt{enable\_if}. The standard's \texttt{std::advance} dispatches on iterator category:

\begin{lstlisting}[language=C++, caption={tag dispatch on iterator category}]
namespace detail {
    template <class It, class D>
    void advance(It& it, D n, std::random_access_iterator_tag) { it += n; }  // O(1)
    template <class It, class D>
    void advance(It& it, D n, std::bidirectional_iterator_tag) {             // O(n)
        for (; n > 0; --n) ++it;
        for (; n < 0; ++n) --it;
    }
    template <class It, class D>
    void advance(It& it, D n, std::input_iterator_tag) { for(; n>0; --n) ++it; }
}
template <class It, class D>
void advance(It& it, D n) {
    detail::advance(it, n,
        typename std::iterator_traits<It>::iterator_category{}); // pick by tag
}
\end{lstlisting}

The tag is an empty struct the optimizer erases; its sole purpose is to steer overload resolution.

%% ============================================================
\section{Type Aliases and Alias Templates (\texttt{using})}

C++11's \texttt{using} replaces \texttt{typedef} and, crucially, supports \textbf{templated aliases} --- impossible with \texttt{typedef}.

\begin{lstlisting}[language=C++, caption={alias and alias template}]
using Bytes = unsigned char;                 // clearer left-to-right than typedef
template<typename T>
using Dictionary = std::map<std::string, T>; // alias TEMPLATE
Dictionary<int> scores;                      // std::map<std::string, int>
\end{lstlisting}

Alias templates are the idiomatic way to provide trait shortcuts (the \texttt{\_t}/\texttt{\_v} convention). The standard added \texttt{\_t} aliases in C++14 and \texttt{\_v} variable templates in C++17; in C++11 you write the \texttt{typename ...::type} form or define the aliases yourself.

%% ============================================================
\section{\texttt{std::tuple}}

\texttt{std::tuple<Ts...>} generalizes \texttt{std::pair} to any number of heterogeneous elements --- a variadic product type.

\begin{lstlisting}[language=C++, caption={tuple basics}]
#include <tuple>
std::tuple<int, double, std::string> t(1, 3.14, "hi");
int    i = std::get<0>(t);            // access by index
auto   u = std::make_tuple(42, 'x');  // type-deduced construction
std::size_t n = std::tuple_size<decltype(t)>::value; // 3
\end{lstlisting}

Unpacking a tuple into a function call is the classic variadic exercise: generate a compile-time index sequence, then expand \texttt{std::get<Is>(tuple)...} as the call's arguments. (C++14 standardized \texttt{std::index\_sequence}; C++17 standardized this as \texttt{std::apply}.)

\begin{lstlisting}[language=C++, caption={calling a function with arguments stored in a tuple}]
namespace detail {
    template <class F, class Tuple, std::size_t... Is>
    auto apply_impl(F&& f, Tuple&& t, index_sequence<Is...>)
        -> decltype(std::forward<F>(f)(std::get<Is>(std::forward<Tuple>(t))...)) {
        return std::forward<F>(f)(std::get<Is>(std::forward<Tuple>(t))...);
    }
}
// apply(f, make_tuple(42,'x',3.14)) calls f(42,'x',3.14)
\end{lstlisting}

%% ============================================================
\section{Professional Insights}

\textbf{Recursion depth is bounded.} Compilers cap template instantiation depth (g++ defaults to $\sim$900; older builds 256). Deeply recursive metafunctions can exceed it --- raise it with \texttt{-ftemplate-depth=N}, or restructure to logarithmic depth for large packs.

\textbf{Prefer \texttt{constexpr} over struct metafunctions when you can.} A \texttt{constexpr} function reads and debugs like ordinary code and serves both compile-time and runtime callers from one definition. Reserve struct metafunctions for \emph{type-level} computation that \texttt{constexpr} cannot express.

\textbf{TMP is a build-time cost.} Every distinct instantiation is compiled; heavy trait machinery and large variadic expansions inflate compile times and binary size. Measure, and cache common instantiations behind \texttt{extern template} (Chapter~\ref{ch:advanced-core-language-and-literals}) where appropriate.

\textbf{\texttt{if constexpr} is the future, but not a total replacement.} C++17's \texttt{if constexpr} cleans up \emph{implementation} branching that once required \texttt{enable\_if}, but does not replace SFINAE for \emph{overload-set} and open extensibility, where tag dispatch still wins. C++11 has neither --- recursion, the expander trick, \texttt{enable\_if}, and tag dispatch are your tools.

\textbf{Name your constraints.} A trait named \texttt{is\_sizeable<T>} documents intent far better than an inline \texttt{decltype(...)} leaking implementation into a signature.
